\documentclass[12pt]{letter}
\usepackage[a4paper,margin=2.5cm]{geometry}
\usepackage{times}
\usepackage{hyperref}

\signature{%
  Zhiqiang Li\\
  College of Physics\\
  Sichuan University\\
  Chengdu 610064, China\\
  \texttt{zhiqiangli@scu.edu.cn}
}

\address{%
  College of Physics\\
  Sichuan University\\
  Chengdu 610064, China
}

\date{\today}

\begin{document}

\begin{letter}{%
  Editor-in-Chief\\
  \emph{Energy Conversion and Management}\\
  Elsevier
}

\opening{Dear Editor,}

We are pleased to submit our manuscript entitled
\textbf{``A Unified Continuous-Time Electrochemical--Power Coupled Framework
for Portable Energy Storage Management and Time-to-Empty Prediction''}
for consideration for publication in \emph{Energy Conversion and Management}.

\medskip
\noindent\textbf{Scope and relevance.}\quad
The manuscript presents a physics-based, continuous-time framework that
couples a modified Shepherd electrochemical model with component-resolved
power profiles to predict the time-to-empty (TTE) of lithium-ion cells
in portable energy storage systems.  Unlike conventional coulomb-counting
or data-driven approaches, our model is formulated entirely in closed-form
ordinary differential equations (ODEs), enabling real-time execution on
low-power microcontrollers (ARM Cortex-M class, $<$\,1\,kB RAM).
We believe this work falls squarely within the scope of \emph{Energy
Conversion and Management} because it addresses the \emph{system-level
energy conversion chain}---from electrochemical state to workload-adaptive
power delivery---and proposes an \emph{optimised energy management strategy}
that demonstrably extends device runtime.

\medskip
\noindent\textbf{Key contributions.}\quad
\begin{enumerate}
  \item A unified ODE framework merging a modified Shepherd battery model
        (RMSE 17.85\,$\pm$\,0.44\,mV on the eight-cell fitting set, with
        cross-batch generalisation confirmed on 15 additional cells) with component-level
        power demand, validated on the publicly available XJTU lithium-ion
        battery dataset.
  \item An adaptive energy management strategy delivering an estimated
        6.1\,Wh energy saving per charge cycle ($\approx$\,2.2\,kWh per
        device per year), outperforming commercial heuristics such as
        Android Battery Saver and iOS Low Power Mode.
  \item Comprehensive robustness analysis including cross-batch validation,
        bootstrap confidence intervals, ablation studies, and parameter
        sensitivity analysis, demonstrating model generalisability.
  \item A semi-empirical cycle-ageing extension preserving the ODE
        structure, enabling lifecycle energy planning.
\end{enumerate}

\medskip
\noindent\textbf{Novelty statement.}\quad
To the best of our knowledge, no prior work has combined a closed-form
electrochemical battery model with component-resolved power coupling in
a single ODE system that can be solved in real time for TTE prediction.
Our framework bridges the gap between high-fidelity battery modelling
and practical energy management on resource-constrained embedded platforms.

\medskip
\noindent\textbf{Ethical and authorship statements.}\quad
This manuscript has not been published previously and is not under
consideration elsewhere.  All authors have approved the manuscript and
agree with its submission to \emph{Energy Conversion and Management}.
No conflicts of interest exist.

\medskip
\noindent\textbf{Suggested reviewers.}\quad
We respectfully suggest the following independent reviewers with relevant
expertise:
\begin{enumerate}
  \item Prof.\ Kailong Liu, University of Warwick, UK ---
        battery charging management, degradation-aware optimisation
        (\texttt{kailong.liu@warwick.ac.uk})
  \item Prof.\ Rui Xiong, Beijing Institute of Technology, China ---
        battery state estimation, BMS
        (\texttt{rxiong@bit.edu.cn})
  \item Prof.\ Xiaosong Hu, Chongqing University, China ---
        equivalent-circuit modelling, comparative battery studies
        (\texttt{xiaosonghu@ieee.org})
\end{enumerate}

\closing{Yours sincerely,}

\end{letter}
\end{document}
